\documentclass[11pt]{article}
\usepackage{amsmath,amssymb,amsthm,hyperref}
\newtheorem{theorem}{Theorem}
\newtheorem{lemma}[theorem]{Lemma}
\newtheorem{corollary}[theorem]{Corollary}
\newtheorem{proposition}[theorem]{Proposition}
\title{Laplace Sign Consistency and the Cost of Ideal Moment Constructions}
\author{}
\date{November 23, 2026}
\begin{document}
\maketitle

\paragraph{Scope.}
These are obstructions to specified construction mechanisms, not a
lower bound ruling out the requested head complexity. They apply even
when all affine forms depend on the target and the common root multiset
varies with the input. The final explicit example explains why the
restriction is substantive. The proofs are manual, not Lean checked.

\begin{theorem}[Counting affine-weighted exponential scores on the cube]
Let $N=2^n$, $H\ge1$, and $p=(2n+1)H+1\le N$. The number of
strict sign tables of
\[
 E(x)=c+\sum_{h=1}^H A_h(x)\exp(-B_h(x)),
 \qquad x\in\{0,1\}^n,
\tag{1}
\]
where $A_h,B_h$ are arbitrary real affine forms, is at most
\[
 \left(\frac{4e(n+1)N}{p}\right)^p.
\tag{2}
\]
For every fixed $C>0$, the union over $H\le C N/n^2$ therefore
has only $2^{o(N)}$ strict sign tables.
\end{theorem}
\begin{proof}
Write $B_h=b_{h0}+\sum_i b_{hi}x_i$ and put
$\theta_{hi}=\exp(-b_{hi})$. Absorb the positive factor
$\exp(-b_{h0})$ into all coefficients of $A_h$. Since the inputs
are Boolean, (1) becomes
\[
 c+\sum_h\left(a_{h0}+\sum_i a_{hi}x_i\right)
                         \prod_{i:x_i=1}\theta_{hi}.
\tag{3}
\]
At each fixed input this is a polynomial in the $p$ parameters
$c,a_{hi},\theta_{hi}$, of total degree at most $n+1$.
Enlarging the domain $\theta_{hi}>0$ to all real values only
increases the number of sign tables. Warren's strict-sign theorem
gives (2). We use its formulation in Goldberg--Jerrum,
\emph{Bounding the Vapnik--Chervonenkis Dimension of Concept Classes
Parameterized by Real Numbers},
\href{https://www.cs.ox.ac.uk/people/paul.goldberg/papers/GJ-ML95.pdf}
{Machine Learning 18 (1995), Theorem 2.1, page 137}.
If $p=N$, the trivial bound $2^N$ is already smaller than (2),
so the possible strict inequality convention in that theorem is immaterial.

For $H\le C N/n^2$, one has $p=O_C(N/n)$. The logarithm of
(2) is $O_C(N\log n/n)=o(N)$, uniformly over these $H$.
For sufficiently large $n$, this follows also by monotonicity of
$p\log(4e(n+1)N/p)$ throughout the allowed range. Taking the
union adds at most $O_C(n)$ to the logarithm.
\end{proof}

\begin{corollary}[A shared Laplace time is too restrictive]
Suppose an admissible rational score is
\[
 S(x)=c+\sum_h\frac{A_h(x)}{B_h(x)},\qquad B_h(x)>0,
\]
and let
\[
 E_t(x)=c e^{-t}+\sum_h A_h(x)e^{-tB_h(x)}.
\tag{4}
\]
Then $S(x)=\int_0^\infty E_t(x)\,dt$. At the budget
$H\le C2^n/n^2$, only $2^{o(2^n)}$ target tables can have a
representation for which some common $t>0$ satisfies
$y(x)E_t(x)>0$ at every input. In particular this holds if the
inequality is required for all positive times.
\end{corollary}
\begin{proof}
Integrate the finite sum in (4), using $B_h(x)>0$.
For any selected common $t$, its multiplication of every denominator
coefficient can be absorbed into the freely chosen affine form $B_h$,
and $ce^{-t}$ is simply a new bias. Thus even a target-dependent
choice of $t$ gives a table in (1). Apply the theorem.
\end{proof}

\begin{lemma}[A positive divided exponential sum]
For distinct real numbers $b_1,\ldots,b_m$ and every $t>0$,
\[
 K_b(t)=\sum_{h=1}^m
       \frac{e^{-tb_h}}{\prod_{j\ne h}(b_j-b_h)}>0.
\tag{5}
\]
\end{lemma}
\begin{proof}
For $m=1$ this is immediate. Otherwise let $q$ be the degree-at-most
$m-1$ Lagrange interpolant to $g(z)=e^{-tz}$ at the $m$ nodes.
Its leading coefficient is
$\sum_h g(b_h)/\prod_{j\ne h}(b_h-b_j)$.
Repeated Rolle's theorem applied to $g-q$ shows that this coefficient
equals $g^{(m-1)}(\xi)/(m-1)!$ for some point between the smallest
and largest node. Multiplication by $(-1)^{m-1}$ gives
$K_b(t)=t^{m-1}e^{-t\xi}/(m-1)!>0$.
\end{proof}

\begin{theorem}[All ideal first variations have consistent Laplace signs]
Let $u_h,v_h,a_h,b_h$ be affine forms, $1\le h\le m$.
At every Boolean input, suppose $u$ and $v$ have the same multiset
of values, with no repetitions. This multiset may depend on the input.
Suppose the product first variation
\[
 R_x(s)=\left.\partial_\epsilon\right|_0
 \left[\prod_h(s+u_h(x)+\epsilon a_h(x))-
       \prod_h(s+v_h(x)+\epsilon b_h(x))\right]=r(x)
\tag{6}
\]
is independent of $s$ and has strict signs $y(x)$.
Then for every $t>0$,
\[
 y(x)\left(\sum_h a_h(x)e^{-tu_h(x)}
             -\sum_h b_h(x)e^{-tv_h(x)}\right)>0
\tag{7}
\]
at every Boolean input. Therefore this entire ideal-first-variation
construction, including moving root multisets, is not universal with
$m=O(2^n/n^2)$.
\end{theorem}
\begin{proof}
Fix an input and match the two lists by their distinct values.
Write $\pi(h)$ for the index with $v_{\pi(h)}=u_h$ and
$d_h=\prod_{j\ne h}(u_j-u_h)$. Evaluation of (6) at $s=-u_h$
gives
\[
 a_h-b_{\pi(h)}=r/d_h.
\]
Consequently the expression in parentheses in (7) is exactly
$rK_u(t)$, positive after multiplication by $y$ by (5).
At any one fixed $t$, (7) is a $2m$-term score in (1), with
affine weights. Apply the counting theorem. No assumption that
the roots are constant across inputs was used.

The distinctness condition loses no nonzero ideal first variation.
If a base value is repeated, each summand in either product's first
derivative still vanishes at the negative of that value. Evaluation
of (6) there would force $r=0$ at that input.
\end{proof}

\begin{lemma}[Positive convolution and its elementary Laplace transform]
For $b_1,\ldots,b_k>0$, let $k_b(t)=e^{-bt}$ on $t\ge0$.
The iterated convolution $k_{b_1}*\cdots*k_{b_k}$ is strictly
positive for every $t>0$ and has Laplace transform
$\prod_{j=1}^k(s+b_j)^{-1}$ for $s\ge0$.
These assertions permit repeated $b_j$.
\end{lemma}
\begin{proof}
Positivity follows by induction from the integral over $0<u<t$.
The transform identity follows by the substitution $t=u+v$ in the
integral of a convolution; the nonnegative integrand justifies
interchanging the integrals. The transform of one factor is
$1/(s+b)$.

We also record the uniqueness fact needed below without requiring a
general inversion theorem. Finite convolutions of these exponentials
are finite sums of $t^j e^{-bt}$: this follows inductively by
integrating an exponential times a polynomial, with equal exponents
giving one extra power of $t$. The Laplace transform of
$t^j e^{-bt}$ is $j!/(s+b)^{j+1}$, by integration by parts.
If a finite sum of these transforms vanishes, uniqueness of partial
fractions forces every coefficient to vanish. That uniqueness follows
by multiplying by the highest power at a pole, evaluating there,
and then successively removing the remaining powers. Hence equal
transforms in this finite class imply equal functions.
\end{proof}

\begin{theorem}[Boolean ideal product differences are also excluded]
Let $u_1,\ldots,u_m,v_1,\ldots,v_m$ be real affine forms and suppose,
at every Boolean input, that
\[
 F(s,x)=\prod_h(s+u_h(x)),\quad
 G(s,x)=\prod_h(s+v_h(x)),\quad F(s,x)-G(s,x)=r(x),
\tag{8}
\]
with $r$ independent of $s$ and strictly of sign $y(x)$.
These identities need only hold on the cube, not as polynomial
identities in the input coordinates. Then
\[
 y(x)\left(\sum_h e^{-tv_h(x)}-\sum_h e^{-tu_h(x)}\right)>0
                  \qquad(t>0).
\tag{9}
\]
Thus ideal product differences and their logarithmic-derivative head
compiler cannot give the requested universal bound at
$m=O(2^n/n^2)$.
\end{theorem}
\begin{proof}
Choose a common affine shift $C(x)$ so that all coefficients of
$C+u_h,C+v_h$ are strictly positive. This is possible by taking
each coefficient of $C$ larger than one plus the largest absolute
coefficient in the two lists. Fix an input and temporarily replace
the roots by their shifted positive values. Substitution
$s\mapsto s+C(x)$ leaves $F-G=r$ unchanged.

Since $F_s=G_s$, logarithmic differentiation gives
\[
 \frac{G_s}{G}-\frac{F_s}{F}
   =\frac{r F_s}{FG}
   =r\sum_{j=1}^m\frac{1}{s+u_j}
                          \prod_{\ell=1}^m\frac{1}{s+v_\ell}.
\tag{10}
\]
The left side is the Laplace transform of
$\sum_h e^{-tv_h}-\sum_h e^{-tu_h}$. By the lemma, the
right side is the transform of $r$ times a sum of $m$ strictly
positive convolution kernels. The uniqueness fact in the lemma
identifies the two functions. This proves (9) for the shifted roots;
removing the common positive factor $e^{-tC(x)}$ proves it as stated.
Finally apply (2) at $t=1$, even allowing affine rather than
constant numerator weights. Repeated roots cause no difficulty here.
\end{proof}

\begin{theorem}[Ideal leading terms at arbitrary analytic parameter order]
Let $u_h(\epsilon,x),v_h(\epsilon,x)$, $1\le h\le m$, be real
analytic in $\epsilon$ near zero, with affine dependence on $x$
for each $\epsilon$. Suppose their products satisfy, at every Boolean
input, the coefficientwise expansion in the auxiliary variable $s$
\[
 F_\epsilon(s,x)=\prod_h(s+u_h(\epsilon,x)),\qquad
 G_\epsilon(s,x)=\prod_h(s+v_h(\epsilon,x)),
\]
\[
 F_\epsilon(s,x)-G_\epsilon(s,x)
    =\epsilon^k r(x)+O(\epsilon^{k+1}),\qquad k\ge1,
\tag{11}
\]
where $r(x)$ is independent of $s$ and has strict signs $y(x)$.
Then $y$ has an affine-weighted exponential representation with at
most $2m$ terms. The base root multiset may depend on $x$ and may
have repetitions. Thus increasing the order or coalescing base roots
in this ideal-leading-term strategy does not overcome (2).
\end{theorem}
\begin{proof}
The order-zero products agree; write their common value as $P(s,x)$.
After a common affine shift, all base root values are positive.
They remain at least some $\delta>0$ for sufficiently small
$|\epsilon|$, uniformly over the finite cube and finite root lists.
Fix one input for the intermediate calculation. The logarithmic
derivative difference has expansion
\[
 \frac{(G_\epsilon)_s}{G_\epsilon}
       -\frac{(F_\epsilon)_s}{F_\epsilon}
       =\epsilon^k r\frac{P_s}{P^2}
                                 +O(\epsilon^{k+1}).
\tag{12}
\]
Indeed its numerator is
$(G_\epsilon)_s(F_\epsilon-G_\epsilon)
 -G_\epsilon((F_\epsilon)_s-(G_\epsilon)_s)$.
In (11) the derivative in $s$ of the order-$k$ coefficient is zero.
Taking the coefficient of $\epsilon^k$ proves (12); all smaller
coefficients vanish.

The left side of (12) is the Laplace transform of the finite sum
\[
 T_\epsilon(t,x)=\sum_h e^{-t v_h(\epsilon,x)}
                         -\sum_h e^{-t u_h(\epsilon,x)}.
\]
Coefficients through order $k$ may be taken under this integral.
Here is a domination justification. By induction in $j$, the
$j$th parameter derivative of $e^{-t u_h}$ is $e^{-t u_h}$
times a polynomial in $t$ of degree at most $j$: differentiate
the polynomial and subtract $t(\partial_\epsilon u_h)$ times
it to pass from $j$ to $j+1$. On a sufficiently small closed
parameter interval, all coefficients of these finitely many
polynomials are bounded. Thus the derivatives through order $k$
are bounded by $M(1+t)^k e^{-\delta t}$ for some $M$, which
is integrable on $[0,\infty)$. The same argument applies to $v_h$.

Every parameter coefficient of $T_\epsilon$ is a finite exponential
polynomial in $t$. The transform uniqueness in the preceding lemma
therefore makes every coefficient below order $k$ zero. At order
$k$ it identifies the coefficient with the inverse transform of
$rP_s/P^2$. If $P=\prod_j(s+b_j)$, then
\[
 \frac{P_s}{P^2}=\sum_j\frac1{s+b_j}\prod_\ell\frac1{s+b_\ell}.
\]
Its inverse transform is a sum of strictly positive convolution
kernels for every $t>0$, including when some $b_j$ coincide.
Consequently the leading coefficient of $T_\epsilon(1,x)$ has
strict sign $y(x)$. On the finite cube, one sufficiently small
$\epsilon>0$ makes every $T_\epsilon(1,x)$ have that sign.
Its $2m$ exponents are affine in $x$ by hypothesis, so it belongs
to (1), even with constant numerator weights. Removing the common
shift only multiplies the exponential sum by a positive factor.
Apply (2).
\end{proof}

\begin{corollary}[Uniform-sign Mellin lifts cannot compress enough]
The uniform-sign Laurent-polynomial construction in
\texttt{toric-switching-audit.tex} is not universal at
$H=O(2^n/n^2)$ monomials, regardless of how large its exponents are.
\end{corollary}
\begin{proof}
Evaluate its asserted uniform sign at $t=1/2$.
Each monomial becomes a constant times
$\prod_{i:x_i=1}\theta_{hi}$, with positive $\theta_{hi}$.
This is the constant-weight subfamily of (3). Arbitrarily large
integer or real exponents are already included by allowing arbitrary
positive $\theta_{hi}$, so they do not invalidate the count.
\end{proof}

\begin{proposition}[General admissible scores can escape the restriction]
On the one-bit cube the admissible two-head score
\[
 S(x)=\frac{1+x}{1+9x}+\frac{-2+x}{3+3x}
\]
is strictly positive at both vertices. Nevertheless its exponential
integrands are never simultaneously strictly positive at any
common $t>0$.
\end{proposition}
\begin{proof}
The score values are $S(0)=1/3$ and $S(1)=1/30$.
All denominator coefficients are strictly positive. The integrands are
\[
 E_t(0)=e^{-t}-2e^{-3t},\qquad
 E_t(1)=2e^{-10t}-e^{-6t}.
\]
The first is positive exactly when $t>(\log 2)/2$; the second
is positive exactly when $t<(\log 2)/4$. These intervals are
disjoint. Thus the argument above is not a lower bound for the
unrestricted head model. This example is a boundary check, not an
efficient universal construction.
\end{proof}

\paragraph{Research consequence.}
The fixed-root face invariant and the moving-root computational
examples remain valid. The stronger exponential factorization now
closes ideal first variations, ideal product differences, and analytic
ideal leading terms at the target budget, as well as uniform-sign Mellin lifts.
Any proposed universal proof based on these integrals must therefore
use a materially different sign-cancellation mechanism. No such
universal mechanism is proved in this note.

Sums of blocks whose Laplace signs conflict at some inputs are not
excluded by these statements. Neither are leading coefficients with
nonconstant dependence on $s$ whose inverse transforms change sign.
The general head model permits both behaviors.
\end{document}
